% commands.tex — кастомные макросы для книги о платёжном стеке

%% ── Платёжные термины ────────────────────────────────────────────────────────
% Единообразное форматирование специфических терминов по всей книге.

\newcommand{\iso}[1]{\textsc{iso\,#1}}           % \iso{8583} → ISO 8583
\newcommand{\mti}[1]{\texttt{#1}}                 % MTI значение
\newcommand{\de}[1]{DE\,\texttt{#1}}              % Data Element: \de{55}
\newcommand{\rc}[1]{\texttt{#1}}                  % Response Code: \rc{05}
\newcommand{\pan}[1]{\texttt{#1}}                 % PAN/номер карты
\newcommand{\hex}[1]{\texttt{0x#1}}               % hex-значение
\newcommand{\law}[1]{\textit{#1}}                 % нормативный акт: \law{161-ФЗ}
\newcommand{\scheme}[1]{\textsc{#1}}              % платёжная схема: \scheme{Visa}

%% ── Специальные блоки ────────────────────────────────────────────────────────
% Определения, предупреждения и справочные блоки через kaobook environments.

% "На практике" — практический совет
\newenvironment{practice}{%
  \begin{kaobox}[title=На практике]
}{%
  \end{kaobox}
}

% "Важно" — критическое замечание
\newenvironment{important}{%
  \begin{kaobox}[title=Важно, colback=red!5, colbacktitle=red!20]
}{%
  \end{kaobox}
}

% "Определение" — глоссарийная вставка внутри текста
\newenvironment{termdef}[1]{%
  \begin{kaobox}[title={Определение: #1}, colback=blue!4, colbacktitle=blue!15]
}{%
  \end{kaobox}
}

%% ── Аббревиатуры (разворачиваются при первом упоминании) ─────────────────────
% Используйте \gls{pan} вместо "PAN" — пакет glossaries из kaobook.
% Файл с определениями: assets/glossary.tex (добавить при необходимости).

%% ── Todo-заметки ─────────────────────────────────────────────────────────────
% Цветовые категории для \todo{}:
\newcommand{\todofig}[1]{\todo[inline, color=blue!15]{Рисунок: #1}}
\newcommand{\todocite}[1]{\todo[inline, color=green!30]{Источник: #1}}
\newcommand{\todocheck}[1]{\todo[inline, color=orange!40]{Проверить: #1}}

%% ── Вспомогательные команды ──────────────────────────────────────────────────
\newcommand{\tbd}{\textcolor{gray}{\textit{[раздел в работе]}}}
\renewcommand{\eg}{т.\,е.\ }
\newcommand{\cf}{ср.\ }
